% ---
% id: meat-core/rem:optol-direction
% kind: rem
% title: "Scope of the perturbation gap"
% summary: >-
%   The gap of [[prop:optol-nobound]] involves no defect, lives in the
%   acceptance-test penalty and can point either way. Among families of finite
%   value the instance leaves an allowance of 53·α/(α−1) bits per round or more
%   unrefuted; at a key rate below one bit no key-length claim can rest on it.
% prominence: standard
% status: accepted
% uses: [def:feasible, prop:value-range]
% mentions: [prop:optol-nobound, def:max-key-length, rem:alpha-domain, rem:representability-ceiling, thm:keylength, thm:reduction, req:infinite-value]
% assumes: []
% refs: []
% provenance: newly-derived
% lean:
% ---

Proposition~\ref{prop:optol-nobound} perturbs only the measurement
$\{F_v\}$ on $\Qsys$, not $\tau_A$, $\{M_x\}$ or the constraint data, and
it compares values of $\hlow$, the number Theorem~\ref{thm:keylength}
turns into a key length and Theorem~\ref{thm:reduction} places below
$\hqkd$, not values of $\hqkd$ itself. Its pair stops at
$\frac{\alpha}{\alpha-1}\cdot52$, while the figure $53$ in the remark is
the largest value of $\log_2\frac1{g_0}$ at an exactly admissible family
of the instance with $g_0>0$, since such a family has $g_0\ge2^{-53}$
(Remark~\ref{rem:representability-ceiling}); between the two scales the
proposition decides nothing. Its part~(iii) applies when the smaller of
the two key lengths is positive, which depends on $n$ and on the other
terms of $\ell_{\max}$; the proposition does not say for which $n$ that
happens on the instance.

\begin{remark}\label{rem:optol-direction}
Proposition~\ref{prop:optol-nobound} does not say that a small
perturbation always matters: it says that no bound in the declared
metadata alone can say that it does not. It does not locate the
fragility in any defect (a failure of Hermiticity, positivity or
completeness in the declared numbers): both families are exact POVMs, so the phenomenon belongs to the map
$\{F_v\}\mapsto\hlow$ and would be there for a \tref{producer} that emitted
perfect numbers for a slightly different device. It does not depend on
the entropy term: the gap lives entirely in the acceptance-test penalty
$\frac{\alpha}{\alpha-1}\sfD(\mathbf q\|\Phimap[\omega])$ of
\eqref{eq:objective}, with $\mathbf q$ ranging over the acceptance set, a
term every declaration carries, and it is the same unboundedness that
Proposition~\ref{prop:value-range}\,(iii) records for $\hlow$ itself, there
as a range statement and here as a stability statement.

And it does not refute every allowance: an $\mathrm{err}$ whose value on
the data of part~(i) is at least $\frac{\alpha}{\alpha-1}\cdot53$ is not
refuted at any admissible $\epsilon$ by a pair of families of the instance
of finite value, that is with $g_0>0$, and this module neither exhibits nor
rules out one. The family $(0,1)$ is exactly admissible as well and lies
within $2^{-53}\le\epsilon$ of $(2^{-53},1-2^{-53})$, but its value is
$\hlow=+\infty$ (Proposition~\ref{prop:value-range}\,(iv)), which carries
no key-length claim (Requirement~\ref{req:infinite-value}); it is set aside
here, since a gap to an infinite value is not a gap between two rates. On those data such an $\mathrm{err}$
would license a discrepancy of more than $53$ bits per round
($\frac{\alpha}{\alpha-1}>1$ at every $\alpha>1$, and it diverges as
$\alpha\downarrow1$) against a key rate whose entropy part is at most
$\log_2|\Zset|=1$ bit per round there ($d_Z=1$), by the entropy-part
ceiling of Proposition~\ref{prop:value-range}\,(ii) (the bound
$\log_2|\Zset|=\log_2(d_Z+1)$ on each entropy term, the rest of $\hlow$
being the acceptance-test penalty). For the protocols this module rates, a
key rate lies below one bit per round in practice, and against such a rate
a robustness allowance of tens of bits per round is not a statement any
key-length claim can rest on. At
a larger $d_Z$ the entropy-part ceiling grows, and the proposition
constrains $\mathrm{err}$ on the data of part~(i) only.

What the proposition does say is, first, that the direction is not
controlled: the gap \eqref{eq:optol-gap} has a sign, and exchanging the
roles of the two families reverses it, so a perturbation of a declared
family may move $\hlow$ up (for a key-length claim, the unsafe direction)
as easily as down. Second, that no threshold escapes the difficulty by
being small: $\epsilon$ may be taken as small as the number model
permits, the gap being bought by the second parameter of the instance,
$N$, rather than by the size of the perturbation. Third, for
$\alpha\in(1,2)$, that the gap reaches the output: by part~(iii) the key
lengths themselves differ, by at least $nB-1$ bits whenever the smaller is
positive, so the gap is not confined to an unused intermediate quantity. At $\alpha\ge2$ the same computation compares
two values of the expression of Definition~\ref{def:max-key-length}, to which this
module gives no key-length reading (Remark~\ref{rem:alpha-domain}).
\end{remark}
